\section{总结与延伸}

\subsection{一条贯穿全课的主线：先设计 memory interface}

这堂课可以压缩为四步推理。第一，每个 autoregressive model 都有跨步 state。第二，state 的形式决定访问接口：token database 支持精确检索，finite state 支持持续压缩。第三，数据单位的 resolution 决定哪种接口自然；attention 需要有效 token，SSM 更能处理需要先压缩的高分辨率流。第四，H-Net 展示一种组合方式：低层学习 chunk，高层对 chunk 做更强关系建模。

\begin{importantbox}{四句带走}
\begin{enumerate}
  \item Mamba 的核心是 state expansion、selectivity 与 efficient scan 的系统组合。
  \item Attention 的真正资产是 token-resolution random access，不只是一个 $QK^\top$ 公式。
  \item Tokenization / patchification 改变建模问题，而不仅是缩短序列。
  \item 最有前景的系统通常是 resolution-aware hybrid，而不是单一 primitive 通吃所有层。
\end{enumerate}
\end{importantbox}

\subsection{三条证据边界}

第一，perplexity、BPB 与 pretraining loss 是局部指标，不能直接替代 retrieval、reasoning、control 与 downstream science。第二，byte / DNA 结果说明 architecture 与 token resolution 相互作用，不能推出所有原始模态都由 SSM 获胜。第三，H-Net 的 learned boundaries 与 outer-stage ablation 支持 compression hypothesis，却没有唯一识别“有限状态”就是增益因果来源。

\subsection{面向系统设计的检查表}

\begin{enumerate}
  \item \textbf{定义状态：}每个生成步后保存什么？状态是否随 context 增长？
  \item \textbf{定义检索：}任务是否需要精确回看某个历史单位，还是只需 sufficient summary？
  \item \textbf{检查粒度：}当前 token / patch / event 是否有独立语义？边界是否应当 context-dependent？
  \item \textbf{匹配算力：}理论复杂度、HBM traffic、kernel utilization 与 batch behavior 是否一致？
  \item \textbf{做 ablation：}匹配参数、FLOPs、数据和训练 recipe 后，再比较 state type 与 layer ratio。
  \item \textbf{保留边界：}不要把一个 pretraining curve 夸大为普遍智能结论。
\end{enumerate}

\subsection{开放问题}

未来值得追踪的问题包括：dynamic chunking 能否扩展到 frontier scale；hierarchical memory 能否同时保留 coarse summary 与可恢复细节；SSM 的 diffuse dependence 怎样做 causal tracing；learning algorithm 能否摆脱 backprop memory 对超长依赖的限制；hardware 是否会为 true recurrence 提供不同最优点；以及 hybrid 中 attention layer 应由固定周期、动态路由还是任务条件触发。

\subsection{拓展阅读}

\begin{itemize}
  \item Albert Gu，\href{https://goombalab.github.io/blog/2025/tradeoffs/}{On the Tradeoffs of SSMs and Transformers}：本课的长文版概念来源。
  \item Gu 与 Dao，\href{https://arxiv.org/abs/2312.00752}{Mamba}；Dao 与 Gu，\href{https://arxiv.org/abs/2405.21060}{Mamba-2 / Structured State Space Duality}。
  \item Lahoti 等，\href{https://arxiv.org/abs/2603.15569v1}{Mamba-3 v1}：课堂前一个月发布的 inference-first SSM 更新。
  \item Hwang、Wang 与 Gu，\href{https://arxiv.org/abs/2507.07955v2}{H-Net v2}：dynamic chunking 的完整方法与 scaling 结果。
  \item Wang 等，\href{https://arxiv.org/abs/2401.13660}{MambaByte}：byte-level selective SSM；Shah 等，\href{https://arxiv.org/abs/2602.10603v3}{dnaHNet v3}：课堂快照中的 genomic hierarchy。
\end{itemize}

\end{document}
